% Lösungen Einheit 2 von 4 – Winkel zwischen Vektoren, Kanten und Geraden (zusammenbau v0.1)
\einheitenkopf{Einheit 2 von 4 · Winkel zwischen Vektoren, Kanten und Geraden}

\erg{11}{a) Formelwinkel}
\erg{12}{a) $\overrightarrow{AB}$ und $\overrightarrow{AC}$ \quad b) $\overrightarrow{PQ} = (2 | 1 | 2)$, $\overrightarrow{PR} = (3 | -2 | 1)$; $\mathrm{cos}\,\varphi = \frac{6}{3 \cdot \sqrt{14}}$, $\varphi \approx 57{,}7^\circ$ \quad c) $\overrightarrow{DA} = (-1 | -3 | -2)$, $\overrightarrow{DC} = (2 | 1 | 0)$; $\mathrm{cos}\,\varphi = \frac{-5}{\sqrt{14} \cdot \sqrt{5}}$, $\varphi \approx 126{,}7^\circ$ \quad d) $\overrightarrow{QP} = (-4 | -2 | 1)$, $\overrightarrow{QR} = (-3 | 3 | 2)$; $\mathrm{cos}\,\varphi = \frac{8}{\sqrt{21} \cdot \sqrt{22}}$, $\varphi \approx 68{,}1^\circ$ \quad e) $\overrightarrow{BA} \circ \overrightarrow{BC} = (-2 | -2 | -1) \circ (1 | -2 | 2) = 0$: rechter Winkel bei B \quad f) $\mathrm{cos}\,\varphi = \frac{\lvert -4 \rvert}{3 \cdot 3}$, $\varphi \approx 63{,}6^\circ$ \quad g) $\lvert\overrightarrow{AB}\rvert = \lvert\overrightarrow{AC}\rvert = \lvert\overrightarrow{BC}\rvert = \sqrt{8}$, das Dreieck ist gleichseitig; beide Winkel $= 60^\circ$}
\erg{13}{a) Fehler: $\overrightarrow{QP}$ zeigt zum Scheitel P hin, so entsteht der Nebenwinkel; richtig mit $\overrightarrow{PQ} = (2 | 2 | 1)$: $\mathrm{cos}\,\varphi = \frac{6}{3 \cdot \sqrt{17}}$, $\varphi \approx 61{,}0^\circ$ \quad b) Nur dann liegen beide Pfeile auf den Schenkeln des Innenwinkels; zeigt einer zu A hin, ändert das Skalarprodukt sein Vorzeichen, und die Formel liefert den Nebenwinkel \quad c) $\overrightarrow{SP} = (-16 | -11{,}66 | 0)$, $\overrightarrow{SQ} = (-16 | -4{,}34 | 0)$; $\varphi \approx 20{,}9^\circ$}
